% !TeX program = xelatex
\documentclass[11pt,a4paper]{article}
\usepackage{fontspec}
\setmainfont{lmroman10-regular.otf}[
  BoldFont=lmroman10-bold.otf,
  ItalicFont=lmroman10-italic.otf,
  BoldItalicFont=lmroman10-bolditalic.otf
]
\setmonofont{lmmono10-regular.otf}[BoldFont=lmmonolt10-bold.otf]
\usepackage[a4paper,margin=23mm]{geometry}
\usepackage{amsmath,amssymb,booktabs,tabularx}
\usepackage{tikz}
\usetikzlibrary{arrows.meta,positioning}
\usepackage[hidelinks,bookmarksnumbered]{hyperref}
\hypersetup{
  pdftitle={Griductive Solver - Báo cáo đồ án 2},
  pdfsubject={CSC14003 Nhập môn Trí tuệ nhân tạo},
  pdfauthor={Võ Duy Phát; Phan Phước Quang Minh; Bùi Trung Hậu; Danh Thiện Nhân; Nguyễn Châu Tuấn Kiệt}
}
\definecolor{navy}{HTML}{18354A}
\definecolor{mist}{HTML}{EAF1F5}
\renewcommand{\contentsname}{Mục lục}
\renewcommand{\refname}{Tài liệu tham khảo}
\renewcommand{\figurename}{Hình}
\renewcommand{\tablename}{Bảng}
\setcounter{tocdepth}{1}
\setlength{\parskip}{4pt plus 1pt}
\setlength{\parindent}{0pt}
\setlength{\emergencystretch}{2em}
\setlength{\tabcolsep}{5pt}
\renewcommand{\arraystretch}{1.16}
\raggedbottom
\newcolumntype{Y}{>{\raggedright\arraybackslash}X}
\newcolumntype{L}[1]{>{\raggedright\arraybackslash}p{#1}}
\newcommand{\code}[1]{\texttt{\detokenize{#1}}}
\newcommand{\pathcode}[1]{\nolinkurl{#1}}
\newcommand{\SAT}{\texttt{SAT}}
\newcommand{\UNSAT}{\texttt{UNSAT}}
\newcommand{\KB}{\mathrm{KB}}
\newcommand{\studentengine}{Võ Duy Phát}
\newcommand{\studentagent}{Phan Phước Quang Minh}
\newcommand{\studentcnf}{Bùi Trung Hậu}
\newcommand{\studentdpll}{Danh Thiện Nhân}
\newcommand{\studentgui}{Nguyễn Châu Tuấn Kiệt}

\begin{document}
\hypersetup{pageanchor=false}
\begin{titlepage}
\centering
{\large Đại học Quốc gia Thành phố Hồ Chí Minh\par}
{\large\bfseries Trường Đại học Khoa học Tự nhiên\par}
{\large Khoa Công nghệ Thông tin\par}
\vfill
{\large CSC14003 -- Nhập môn Trí tuệ nhân tạo\par}
\vspace{9mm}
{\Large Báo cáo đồ án 2\par}
\vspace{4mm}
{\Huge\bfseries\color{navy} Griductive Solver\par}
\vspace{3mm}
{\Large Suy luận mệnh đề trên sân bóng\par}
\vspace{6mm}
{\color{navy}\rule{65mm}{0.7pt}\par}
\vspace{12mm}
{\large\bfseries Thành viên nhóm\par}
\vspace{3mm}
\begin{tabular}{@{}lr@{}}
\toprule
\textbf{Họ và tên} & \textbf{MSSV}\\
\midrule
\studentengine & 25122013\\
\studentagent & 25122029\\
\studentcnf & 25122006\\
\studentdpll & 25122032\\
\studentgui & 25122026\\
\bottomrule
\end{tabular}
\vfill
Thành phố Hồ Chí Minh, tháng 11 năm 2026
\end{titlepage}
\hypersetup{pageanchor=true}

\section*{Tóm tắt}
Griductive Solver là phiên bản chủ đề bóng đá của trò chơi suy luận
Griductive~\cite{original}. Người chơi chỉ được lật một thẻ khi trạng thái
của cầu thủ được chứng minh từ các manh mối và phán quyết công khai. Nhóm
biểu diễn mỗi cầu thủ bằng một biến Boolean, mã hóa manh mối có cấu
trúc sang dạng chuẩn hội (CNF) và dùng bộ giải SAT bằng DPLL viết
trong dự án để loại trừ phán quyết đối lập. Game Engine quản lý bàn
chơi và dữ liệu ẩn, Controller điều phối lượt chơi, còn Logic Agent
suy luận trên ảnh chụp trạng thái công khai. Giao diện trình duyệt
cục bộ hỗ trợ chơi thủ công, gợi ý và giải tự động. Khi kiểm định
phiên bản \code{b373f86}, cả năm bàn chuẩn từ $3\times3$ đến $5\times5$
đều có nghiệm duy nhất và giải được theo trình tự mở thẻ. Tất cả
25 lượt benchmark đều hoàn thành. Bàn $5\times5$ cần 268 truy vấn SAT
và mất trung bình 6,816 giây cho vòng giải tự động.

\tableofcontents
\clearpage

\section{Kế hoạch và phân công công việc}
Năm thành phần chính của dự án gồm Engine, Agent, Encoder, bộ giải
DPLL và giao diện. Kế hoạch tích hợp đi từ cấu trúc bàn và hợp đồng
dữ liệu đến mã hóa, truy vấn SAT và lượt chơi hoàn chỉnh. Bảng
\ref{tab:team} ghi phạm vi phân công đã có trong hồ sơ nhóm. Cột
bàn giao mô tả các đầu ra theo phân công và vai trò của chúng khi tích
hợp. Tỷ lệ đóng góp cá nhân cần nhóm xác nhận.

\begin{table}[htbp]
\caption{Phân công chính và điểm nối giữa các thành phần.}
\label{tab:team}
\small
\begin{tabularx}{\textwidth}{@{}L{35mm}L{26mm}Y L{21mm}@{}}
\toprule
\textbf{Thành viên} & \textbf{Phụ trách} & \textbf{Đầu ra và vai trò tích hợp} & \textbf{Hoàn thành}\\
\midrule
\studentengine\newline 25122013
& Game Engine
& \pathcode{engine.py}, dữ liệu bàn và quy trình mở thẻ. Cung cấp trạng
thái công khai cho Controller và kiểm tra tính nhất quán với lời giải
khi tiếp nhận phán quyết đã được chứng minh.
& Chờ xác nhận\\
\addlinespace
\studentagent\newline 25122029
& Logic Agent
& \pathcode{agent.py} và điểm tích hợp \pathcode{controller.py}. Ghép
Encoder với DPLL để phân loại, chọn nước đi và ghi nhật ký mở thẻ.
& Chờ xác nhận\\
\addlinespace
\studentcnf\newline 25122006
& CNF Encoding
& \pathcode{cnf.py}, ngữ nghĩa manh mối và phân giải vùng. Chuyển trạng
thái công khai thành các mệnh đề và ánh xạ biến để Agent truy vấn.
& Chờ xác nhận\\
\addlinespace
\studentdpll\newline 25122032
& DPLL SAT Solver
& \pathcode{dpll.py}. Nhận CNF và giả định tạm thời, trả trạng thái
SAT, mô hình khi có nghiệm và các bộ đếm tìm kiếm.
& Chờ xác nhận\\
\addlinespace
\studentgui\newline 25122026
& GUI
& \pathcode{gui.py} và giao diện trình duyệt trong \pathcode{web/}.
Gửi thao tác chơi qua API, dựng bàn và trình bày manh mối, nhật ký,
thống kê. Phân công còn gồm chuẩn bị video minh họa giao diện.
& Chờ xác nhận\\
\bottomrule
\end{tabularx}
\end{table}

Ba điểm nối giúp các phần có thể kiểm thử độc lập. Thứ nhất,
\code{PublicGameState} giới hạn dữ liệu chuyển từ Engine sang Agent.
Thứ hai, một công thức CNF mang mệnh đề và ánh xạ biến, còn kết quả SAT
mang cờ thỏa được, mô hình và thống kê. Thứ ba, Controller nhận kết quả
phân loại để điều phối mở thẻ, sau đó Web API gửi trạng thái mới.
Module \pathcode{contracts.py} định nghĩa các giao diện cấu trúc để
thay Encoder, bộ giải hoặc Engine bằng đối tượng kiểm thử tương thích.
Các kiểm thử ở Mục~\ref{sec:tests} kiểm tra cả hợp đồng lẫn luồng tích
hợp thực tế.

\clearpage
\section{Mô hình bàn chơi và cơ sở tri thức}
Bàn là lưới vuông $N\times N$, mỗi ô có tọa độ, hồ sơ cầu thủ và
một manh mối. JSON của bàn cung cấp lời giải đầy đủ, manh mối theo
chủ thẻ và danh sách thẻ mở lúc đầu. \code{Puzzle.load} đọc JSON
UTF-8 và dựng \code{Puzzle}, \code{CellInfo}, \code{Clue} và
\code{RegionSpec}. Với dữ liệu \code{roster: football}, danh sách
cầu thủ theo cỡ bàn được lấy từ \pathcode{roster.py}. Bàn tùy chỉnh
có thể khai báo ô, tên và nghề nghiệp. Tọa độ đi theo từng hàng từ
trên xuống, trong mỗi hàng đi từ trái sang phải.

Biến chính $C_i$ tương ứng với ô $i$. Trong một mô hình SAT,
$C_i=1$ biểu thị \code{CRIMINAL} và $C_i=0$ biểu thị \code{INNOCENT}.
Trong trạng thái công khai, \code{UNKNOWN} chỉ mức tri thức khi cả hai
giá trị còn khả dĩ. Nó được biểu diễn bằng \code{None} trong
\code{Classification.status} nếu KB nhất quán. Trên giao diện,
\code{UNKNOWN} là nhãn hiển thị cho mọi thẻ chưa mở, kể cả thẻ đã có
thể chứng minh nhưng người chơi chưa thực hiện phán quyết.

\subsection{Manh mối và vùng}
\code{resolve_region} chuyển một \code{RegionSpec} thành danh sách ô
theo tọa độ bàn. Các loại cơ bản gồm hàng, cột, danh sách ô tường
minh và lân cận tối đa tám ô (có đường chéo, bỏ ô trung tâm).
Các vùng mở rộng gồm toàn bàn, góc, biên, các ô nằm giữa hai đầu
mút thẳng hàng hoặc chéo, lân cận chung, ô kề theo một hướng,
nhóm nghề nghiệp, hợp và giao. Bộ phân giải từ chối vùng rỗng
và các tham chiếu sai. Khái niệm kề cạnh dùng trong manh mối liên
thông được xác định riêng ở Mục~\ref{sec:extensions}.

Mỗi manh mối công khai được coi là đúng bất kể trạng thái của chủ thẻ.
Tên và vị trí cầu thủ phục vụ hiển thị, còn thuộc tính nghề nghiệp có thể
tham gia suy luận nếu được manh mối tham chiếu. \code{validate_puzzle}
kiểm tra số ô, thứ tự tọa độ và tên, độ phủ của lời giải và manh mối,
thẻ mở ban đầu, mã định danh manh mối và tham số vùng. Với bàn
thông thường, \code{evaluate_clue} còn đối chiếu từng manh mối với
phép gán ẩn trước khi tạo ván.

\subsection{Trạng thái và giao thức mở thẻ}
Đặt $\mathcal R_t$ là các manh mối đã mở và $P_t$ là các cặp ô,
trạng thái công khai sau $t$ lượt được chấp nhận. Với
$l_i(\mathrm{CRIMINAL})=C_i$ và $l_i(\mathrm{INNOCENT})=\neg C_i$,
\[
 \KB_t =
 \bigwedge_{c\in\mathcal R_t}\operatorname{CNF}(c)
 \ \land\!
 \bigwedge_{(i,s)\in P_t} (l_i(s)).
\]
Bước $t=0$ dùng trạng thái và manh mối của các thẻ mở sẵn.
\code{GameEngine} giữ \code{Puzzle} và dựng \code{PublicGameState}
từ tiêu đề, cỡ bàn, metadata các ô, $\mathcal R_t$ và bản sao $P_t$.
Sau khi Controller chứng minh trạng thái $s$ của ô đang úp $i$,
Engine đối chiếu $s$ với lời giải rồi công bố đúng một thẻ và
manh mối gắn với thẻ đó. Khi ấy $P_{t+1}=P_t\cup\{(i,s)\}$ và
$\mathcal R_{t+1}=\mathcal R_t\cup\{\mathrm{clue}(i)\}$.
Phán quyết bị từ chối giữ nguyên hai tập này.

Một ván hoàn thành khi mọi ô có phán quyết công khai. Vòng tự động
cũng dừng khi không còn phán quyết bắt buộc (\code{STUCK}), khi KB
công khai mâu thuẫn (\code{INCONSISTENT}) hoặc khi chạm giới hạn
số bước mà bên gọi đặt (\code{STEP_LIMIT}).

\clearpage
\section{Kiến trúc và luồng điều khiển}
\label{sec:architecture}
\begin{figure}[htbp]
\centering
\begin{tikzpicture}[
  node distance=8mm and 9mm,
  block/.style={draw=navy!78,fill=mist,rounded corners=2pt,
    align=center,minimum height=9mm,text width=39mm,inner sep=2mm,
    font=\small},
  link/.style={-{Stealth[length=2.1mm]},draw=navy,semithick},
  label/.style={font=\scriptsize,align=center,fill=white,inner sep=1pt}
]
\node[block] (ui) {Trình duyệt\\\code{index.html}, \code{app.js}};
\node[block,right=of ui] (web) {\code{Handler} / \code{Match}\\Web API};
\node[block,right=of web] (launch) {\code{main.py} / \code{gui.py}\\Khởi động server};
\node[block,below=of web] (controller) {\code{GameController}\\Điều phối lượt chơi};
\node[block,below left=12mm and 10mm of controller] (engine) {\code{GameEngine}\\Sở hữu dữ liệu ẩn};
\node[block,below right=12mm and 10mm of controller] (agent) {\code{LogicAgent}\\Phân loại và chọn bước};
\node[block,below=14mm of agent] (dpll) {\code{DPLLSolver}\\SAT / UNSAT, mô hình};
\node[block,left=of dpll] (cnf) {\code{CNFEncoder}\\CNF và ánh xạ biến};
\draw[link,<->] (ui) -- node[label,above]{HTTP} (web);
\draw[link] (launch) -- (web);
\draw[link] (web) -- (controller);
\draw[link] (controller) -| node[label,pos=.28,above]{snapshot, mở thẻ} (engine);
\draw[link] (controller) -| node[label,pos=.43,above]{snapshot công khai} (agent);
\draw[link] ([xshift=-12mm]agent.south) -- ++(0,-6mm)
  -| node[label,pos=.25,above]{mã hóa KB} (cnf.north);
\draw[link] (agent) -- node[label,right]{CNF, giả định} (dpll);
\end{tikzpicture}
\caption{Luồng gọi hàm và đi qua Web API trong một lượt chơi.}
\label{fig:architecture}
\end{figure}

\begin{table}[htbp]
\caption{Dữ liệu vào, kết quả trả về và trách nhiệm của các module.}
\label{tab:interfaces}
\small
\begin{tabularx}{\textwidth}{@{}L{33mm}Y Y@{}}
\toprule
\textbf{Module} & \textbf{Dữ liệu vào} & \textbf{Đầu ra và bên sử dụng}\\
\midrule
\pathcode{main.py}, \pathcode{gui.py}, \pathcode{cli.py}
& Tham số dòng lệnh và đường dẫn JSON.
& Khởi động giao diện cục bộ hoặc chạy \code{solve},
\code{validate}, \code{experiments}.\\
\pathcode{models.py}, \pathcode{roster.py}, \pathcode{regions.py}, \pathcode{clues.py}
& JSON, metadata bàn và khai báo vùng, manh mối.
& Mô hình dữ liệu, danh sách ô, kiểm tra ngữ nghĩa và văn bản mô tả.\\
\pathcode{engine.py}
& Bàn đã phân tích và phán quyết từ Controller.
& Snapshot công khai, manh mối vừa mở và trạng thái hoàn thành.\\
\pathcode{controller.py}
& Ô và phán quyết thủ công hoặc lệnh Hint, Step, Auto Solve.
& Kết quả \code{Submission} hoặc phân loại, trạng thái dừng và trace.\\
\pathcode{agent.py}, \pathcode{contracts.py}
& Snapshot công khai, giao diện Encoder và SAT Solver.
& \code{Classification}, nước đi bắt buộc và thống kê truy vấn.\\
\pathcode{cnf.py}, \pathcode{dpll.py}
& Manh mối và phán quyết đã mở, rồi CNF và giả định tạm.
& Công thức có ánh xạ biến, sau đó là \code{SATResult} và số đo.\\
\pathcode{web/}
& Yêu cầu HTTP và trạng thái do Controller quản lý.
& JSON công khai, tài nguyên tĩnh và giao diện trình duyệt.\\
\bottomrule
\end{tabularx}
\end{table}

CLI và benchmark gọi Engine, Controller hoặc Agent trực tiếp. Trong
luồng chơi, Agent gọi Encoder để tạo công thức và gọi DPLL để giải công
thức đã tạo. Quy trình kiểm định nghiệm duy nhất có luồng gọi riêng với
tập manh mối đầy đủ (Mục~\ref{sec:uniqueness}).

\clearpage
\section{Mã hóa CNF}
\subsection{Biến chính và các manh mối cơ bản}
\code{CNFEncoder} đánh số ô $(r,c)$ bằng $(r-1)N+c$ với hàng và
cột bắt đầu từ 1. Literal là số nguyên có dấu. Chẳng hạn, trên bàn
$3\times3$, \code{B1} có biến số 2 nên $-2$ là phán quyết vô tội.
\code{CNFFormula} lưu riêng \code{primary_variables} và
\code{auxiliary_variables}. Biến phụ bắt đầu tại $N^2+1$ và được tạo
lại cho mỗi công thức, theo thứ tự mã hóa. \code{encode_public_state}
ghép CNF của manh mối công khai và các mệnh đề đơn vị của phán quyết.

Với vùng $R=(i_1,\dots,i_m)$, đặt $\ell_j=C_{i_j}$ nếu đếm
\code{CRIMINAL}, hoặc $\ell_j=\neg C_{i_j}$ nếu đếm
\code{INNOCENT}. Ký hiệu $S_R=\sum_{j=1}^m[\ell_j]$ đếm số literal
đúng. Bảng~\ref{tab:clues} tổng hợp sáu mẫu cơ bản~\cite{spec}.
Mục tiếp theo dẫn xuất các mệnh đề mà bộ mã hóa tự động tạo.

\begin{table}[htbp]
\caption{Ngữ nghĩa và CNF của các mẫu manh mối cơ bản.}
\label{tab:clues}
\small
\begin{tabularx}{\textwidth}{@{}l L{44mm} Y@{}}
\toprule
\textbf{Mẫu} & \textbf{Điều kiện logic} & \textbf{Mệnh đề do Encoder sinh}\\
\midrule
\code{FACT} & Cầu thủ $a$ có trạng thái $s$ & $(l_a(s))$\\
\code{SAME} & $C_a\leftrightarrow C_b$ &
$(\neg C_a\lor C_b)\land(C_a\lor\neg C_b)$\\
\code{DIFFERENT} & $C_a\oplus C_b$ &
$(C_a\lor C_b)\land(\neg C_a\lor\neg C_b)$\\
\code{EXACTLY} & $S_R=k$ & Bộ đếm, chốt $q_{m,k}$.\\
\code{AT_LEAST} & $S_R\ge k$ & Bộ đếm, chốt $\bigvee_{j=k}^{m}q_{m,j}$.\\
\code{AT_MOST} & $S_R\le k$ & Bộ đếm, chốt $\bigvee_{j=0}^{k}q_{m,j}$.\\
\bottomrule
\end{tabularx}
\end{table}

\subsection{Dẫn xuất manh mối nhị phân}
Trong \pathcode{puzzles/3x3_binary_chain.json}, \code{clue_B1}
phát biểu B1 và C1 khác trạng thái. Với $x=C_{\mathrm{B1}}$ và
$y=C_{\mathrm{C1}}$, hai cách để đúng là $(x\land\neg y)$ hoặc
$(\neg x\land y)$. Điều kiện ít nhất một người phạm tội loại
$(0,0)$ bằng $(x\lor y)$. Điều kiện nhiều nhất một người phạm tội
loại $(1,1)$ bằng $(\neg x\lor\neg y)$. Hội hai điều kiện cho
\[
 (x\land\neg y)\lor(\neg x\land y)
 \ \equiv\ (x\lor y)\land(\neg x\lor\neg y).
\]
\code{encode_clue} tạo đúng hai tuple \code{(2, 3)} và
\code{(-2, -3)}. Thẻ mở sẵn A1 có \code{FACT} khẳng định B1
phạm tội. Khi phán quyết ấy mở \code{clue_B1}, đơn vị $(x)$ và
mệnh đề $(\neg x\lor\neg y)$ buộc C1 vô tội.

Với \code{SAME}, hai chiều kéo theo $x\Rightarrow y$ và
$y\Rightarrow x$ lần lượt sinh $(\neg x\lor y)$ và
$(x\lor\neg y)$. Hai mệnh đề giữ lại $(0,0)$ và $(1,1)$. Cả
\code{SAME} lẫn \code{DIFFERENT} đều dùng trực tiếp biến chính.

\clearpage
\subsection{Dẫn xuất manh mối đếm bằng bộ đếm trạng thái}
\label{sec:counter}
\code{_count_states} dùng $q_{i,j}$ để biểu diễn trạng thái đã đọc
$i$ literal và có đúng $j$ literal đúng, với $0\le i,j\le m$.
Ở hàng 0, công thức có $(q_{0,0})$ và $(\neg q_{0,j})$
với $1\le j\le m$. Mỗi hàng $i\ge1$ chứa một mệnh đề buộc có
ít nhất một trạng thái và các mệnh đề loại từng cặp trạng thái
đồng thời
\[
 \bigvee_{j=0}^{m} q_{i,j},
 \qquad
 (\neg q_{i,a}\lor \neg q_{i,b}) \quad (0\le a<b\le m).
\]
Nếu $q_{i-1,j}$ đang đúng, đọc $\ell_i$ sai sẽ giữ nguyên
$j$. Đọc $\ell_i$ đúng sẽ tăng sang $j+1$. Đổi hai phép kéo theo
thành mệnh đề cho
\begin{align*}
 &(q_{i-1,j}\land\neg\ell_i)\Rightarrow q_{i,j}
 &&\leadsto (\neg q_{i-1,j}\lor\ell_i\lor q_{i,j}),\\
 &(q_{i-1,j}\land\ell_i)\Rightarrow q_{i,j+1}
 &&\leadsto (\neg q_{i-1,j}\lor\neg\ell_i\lor q_{i,j+1})
 \quad (j<m).
\end{align*}
Tại $j=m$, nhánh tăng được thay bằng
$(\neg q_{i-1,m}\lor\neg\ell_i)$. Cuối cùng \code{_cardinality}
thêm $\bigvee_{j\in A}q_{m,j}$ với $A$ là tập số đếm được phép.
Ba mẫu \code{EXACTLY}, \code{AT_LEAST}, \code{AT_MOST} lần lượt
dùng $A=\{k\}$, $A=\{k,\ldots,m\}$ và $A=\{0,\ldots,k\}$.
Một điều kiện cho phép mọi số đếm được rút gọn thành CNF rỗng.
Tập số đếm hợp lệ rỗng sinh một mệnh đề rỗng.

\paragraph{Ví dụ từ bàn 3x3 Counting Trail.}
Trong \pathcode{puzzles/3x3_counting_chain.json}, \code{clue_B1}
là \code{EXACTLY(1, ROW 1)}. Đặt $(x,y,z)=(C_{\mathrm{A1}},
C_{\mathrm{B1}},C_{\mathrm{C1}})$. Hai mệnh đề chuyển ở lần đọc thứ
ba từ trạng thái $q_{2,0}$ là
\[
 (\neg q_{2,0}\lor z\lor q_{3,0})
 \ \land\
 (\neg q_{2,0}\lor\neg z\lor q_{3,1}).
\]
Bộ mã hóa sinh $(q_{3,1})$ và tạo đủ các mệnh đề khởi tạo,
chuyển trạng thái và loại trừ ở trên. Riêng manh mối này có
$4^2=16$ biến phụ, 4 mệnh đề khởi tạo, $3(1+6+8)=45$ mệnh đề
trong ba hàng chuyển và 1 mệnh đề chốt, tổng cộng 50 mệnh đề.
Khi mở B1, $\KB_t$ đã chứa $(\neg x)$ và $(\neg y)$. Do đó
$q_{2,0}$ đúng. Nếu giả định $\neg z$, mệnh đề giữ số đếm buộc
$q_{3,0}$, mâu thuẫn với $q_{3,1}$. Vậy C1 phải phạm tội.

Ví dụ ngắn cho hai ngưỡng còn lại nằm ngay trên bàn này.
\code{clue_A1} đếm nhiều nhất 0 trên $\{\mathrm{B1}\}$ nên
chốt $q_{1,0}$, suy ra $\neg C_{\mathrm{B1}}$.
\code{clue_C1} đếm ít nhất 1 trên $\{\mathrm{A2}\}$ nên chốt
$q_{1,1}$, suy ra $C_{\mathrm{A2}}$. Cả hai trường hợp
vẫn đi qua bộ đếm trong mã nguồn.

Với một phép gán các biến chính, quy nạp theo $i$ cho thấy đúng
một $q_{i,j}$ được chọn và $j$ bằng số literal đúng đã đọc. Ngược
lại, nếu phép gán thỏa điều kiện đếm, đặt đúng trạng thái số đếm ở
mỗi hàng cho một mở rộng thỏa CNF. Vì thế công thức có cùng các
phép gán trên biến chính với ngữ nghĩa manh mối. Cài đặt có
$O(m^2)$ biến phụ và $O(m^3)$ mệnh đề do loại trừ theo từng cặp ở
mỗi hàng.

\clearpage
\subsection{Manh mối mở rộng và cách đối chiếu}
\label{sec:extensions}
\begin{table}[htbp]
\caption{Các mở rộng được hỗ trợ trong \texttt{clues.py} và \texttt{cnf.py}.}
\small
\begin{tabularx}{\textwidth}{@{}L{36mm}Y@{}}
\toprule
\textbf{Mẫu} & \textbf{Ngữ nghĩa và cách mã hóa}\\
\midrule
\code{PARITY} & Kiểm tra chẵn/lẻ của số trạng thái đích. Mỗi biến
phụ là XOR của kết quả trước với literal kế tiếp, dùng bốn mệnh đề
cho mỗi phép tương đương. Chốt dấu của kết quả cuối.\\
\code{COUNT_COMPARE} & Dùng hai bộ đếm cho hai vùng. Với mỗi cặp
số đếm $(a,b)$ vi phạm toán tử so sánh, thêm
$(\neg q^{L}_{m_L,a}\lor \neg q^{R}_{m_R,b})$.
Ô thuộc giao hai vùng dùng cùng biến chính ở cả hai bộ đếm.\\
\code{ALL_OF} & Hội của các manh mối con. Ghép danh sách mệnh đề
của từng con. Dữ liệu hiện tại chỉ hỗ trợ một tầng \code{ALL_OF}.\\
\code{INCLUDED_COUNT} & Đếm đúng $k$ trong vùng và buộc cầu thủ
được nêu thuộc trạng thái đích bằng mệnh đề đơn vị.\\
\code{OVERLAP_COUNT} & Ghép một ràng buộc đếm đúng trên vùng
tổng và một ràng buộc đếm đúng trên vùng con.\\
\code{CONNECTED}, \code{CONNECTED_COUNT} & Các ô mang trạng thái đích
liên thông theo bốn hướng. Mẫu thứ hai đồng thời chốt số lượng $k$.\\
\code{GROUP_MOST} & Nhóm được nêu có số đếm lớn hơn nghiêm ngặt
từng nhóm khác. Dùng so sánh \code{GT} cho từng cặp có liên quan.\\
\code{GROUP_ONLY} & Nhóm được nêu có số đếm đúng $k$. Các nhóm
khác lần lượt phải có số đếm thuộc tập $\{0,\ldots,m\}\setminus\{k\}$
với $m$ là kích thước nhóm đó.\\
\bottomrule
\end{tabularx}
\end{table}

\paragraph{Liên thông.}
\code{_encode_connected} buộc có một gốc nếu có ô mang trạng thái
đích. Gốc là ô đích đầu tiên theo thứ tự mà vùng trả về. Từ gốc,
biến đạt tới ở độ sâu 0 đúng tại gốc. Ở độ sâu kế tiếp, một ô
đạt tới nếu đã đạt tới ở bước trước hoặc nếu nó mang trạng thái
đích và có một ô kề cạnh đã đạt tới. Bộ mã hóa tạo biến phụ cho
các hội giữa literal đích và trạng thái đạt tới của hàng xóm, rồi
mã hóa tương đương với hội/tuyển vừa nêu. Sau $m-1$ lớp, mọi ô
đích phải đạt tới. Tập ô đích rỗng thỏa liên thông theo quy ước
của \code{is_connected}. Với \code{GROUP_MOST} và
\code{GROUP_ONLY}, nhóm có thể là hàng, cột, nghề nghiệp hoặc
vùng tám lân cận quanh mỗi người.

\paragraph{Kiểm chứng.}
\code{evaluate_clue} tính trực tiếp ngữ nghĩa trên một phép gán
đầy đủ, độc lập với Encoder. Với phép thử đối chiếu CNF, mỗi
biến chính được cố định bằng một giả định. Bộ giải chỉ cần tìm
phép gán phần biến phụ để kiểm tra tính thỏa được. Các test duyệt
tổ hợp trên bàn nhỏ và thử thêm trường hợp phủ nhau, liên thông,
đếm và xếp hạng trên các bàn lớn hơn. Mục~\ref{sec:tests} ghi số
test thực tế và phạm vi kiểm tra.

\clearpage
\section{Bộ giải SAT bằng DPLL}
\pathcode{dpll.py} cài đặt DPLL theo tìm kiếm nhị phân có lan
truyền đơn vị~\cite{dpll,aima}. Mỗi literal là số nguyên có dấu,
mỗi mệnh đề là tuple literal và CNF là một dãy mệnh đề. Phép
gán bộ phận là \code{dict[int, bool]}. Đầu vào \code{solve}
gồm CNF, chỉ số biến lớn nhất trong công thức sản xuất và dãy
giả định. Giả định được nối vào bản làm việc dưới dạng mệnh đề
đơn vị riêng cho lần gọi. \code{SATResult} chứa cờ thỏa được,
phép gán đầy đủ khi SAT hoặc \code{None} khi UNSAT, và thống kê.

\subsection{Tìm kiếm và lan truyền}
\code{_unit_propagate} quét mệnh đề theo thứ tự. Nếu một mệnh đề
đã có literal đúng, bước quét bỏ qua nó. Nếu chưa thỏa và không
còn literal chưa gán, nhánh hiện tại xung đột. Nếu chỉ còn một
literal chưa gán, biến của literal ấy phải nhận giá trị làm
literal đúng. Bộ giải quét lại tới khi ổn định. Một mệnh đề
rỗng gây xung đột, còn CNF rỗng thỏa ngay.

\code{_search} sao chép phép gán đầu vào trước khi lan truyền.
Khi còn mệnh đề chưa thỏa, nó chọn biến chưa gán có chỉ số nhỏ nhất
trong $1,\ldots,V$ rồi thử \code{False} trước, \code{True} sau.
Đây là lựa chọn theo chỉ số, kể cả biến chưa xuất hiện trong
mệnh đề còn đang xét. Mỗi nhánh nhận bản sao phép gán.
Cài đặt hiện tại quét trực tiếp các mệnh đề và phân nhánh theo
thứ tự cố định.

\begin{figure}[htbp]
\centering
\begin{minipage}{.94\textwidth}
\small
\begin{tabbing}
\hspace{1.6em}\=\hspace{1.6em}\=\hspace{1.6em}\=\kill
\textbf{Solve}$(F,V,A)$\\
\> Tạo thống kê mới và nối mỗi giả định $a\in A$ vào $F$ dưới dạng $(a)$.\\
\> Gọi \textbf{Search}$(F,\{\},V)$ và đo thời gian toàn bộ lần gọi.\\
\> Trả SAT cùng mô hình nếu tìm thấy, ngược lại trả UNSAT.\\[2pt]
\textbf{Search}$(F,M,V)$\\
\> Tăng số lời gọi đệ quy, sao chép $M$.\\
\> Lan truyền đơn vị trên $F$ đến khi ổn định. Nếu có xung đột, trả thất bại.\\
\> Nếu mọi mệnh đề đã thỏa, gán \texttt{False} cho biến còn trống và trả $M$.\\
\> Chọn biến chưa gán nhỏ nhất $v$, tăng số quyết định.\\
\> \textbf{Với} $b$ theo thứ tự \texttt{False}, \texttt{True}\\
\>\> Gọi \textbf{Search}$(F,M\cup\{v\mapsto b\},V)$.\\
\>\> Nếu có mô hình thì trả mô hình. Nếu thất bại thì tăng số quay lui.\\
\> Trả thất bại.
\end{tabbing}
\end{minipage}
\caption{Mã giả tương ứng với \texttt{solve} và \texttt{\_search}.}
\label{fig:dpll-pseudocode}
\end{figure}

\clearpage
\subsection{Thống kê và ví dụ có quay lui}
\begin{table}[htbp]
\caption{Các bộ đếm mà \texttt{DPLLSolver} thu trong mỗi truy vấn.}
\small
\begin{tabularx}{\textwidth}{@{}lY@{}}
\toprule
\textbf{Chỉ số} & \textbf{Sự kiện được ghi nhận}\\
\midrule
\code{decisions} & Một lần chọn biến để rẽ hai nhánh.\\
\code{propagations} & Một phép gán mới vì mệnh đề đơn vị, kể cả do giả định.\\
\code{conflicts} & Một xung đột khiến vòng lan truyền trả thất bại.\\
\code{backtracks} & Một nhánh con thất bại, kể cả nhánh con cuối cùng.\\
\code{recursive_calls} & Mỗi lần vào \code{_search}, kể cả lời gọi gốc.\\
\code{runtime_ms} & Chênh lệch \code{perf_counter} từ đầu \code{solve}, tính bằng ms.\\
\bottomrule
\end{tabularx}
\end{table}

Xét $F=(x\lor y)\land(x\lor\neg y)\land(\neg x\lor z)$ với ID
của $x,y,z$ lần lượt là 1, 2, 3. Gốc chưa có mệnh đề đơn vị
nên chọn $x$. Với $x=0$, mệnh đề đầu buộc $y=1$ còn mệnh đề
thứ hai xung đột. Sau khi quay lui, thử $x=1$. Hai mệnh đề
đầu đã thỏa, mệnh đề cuối buộc $z=1$. Có thể điền $y=0$ vì
mọi mệnh đề đã thỏa. Kết quả là SAT với $(1,0,1)$, một quyết
định, hai lan truyền, một xung đột, một quay lui và ba lần gọi
\code{_search}. Nếu chạy riêng với giả định $-1$, mệnh đề đơn
vị $(\neg x)$ buộc nhánh mâu thuẫn và lần gọi trả UNSAT.

\subsection{Tính đúng đắn, đầy đủ và chi phí}
Lan truyền chỉ gán literal mà mọi phép mở rộng thỏa một mệnh
đề đang đơn vị đều phải gán đúng. Một mệnh đề đã sai hoàn toàn
không thể được cứu bằng cách gán các biến còn lại. Mỗi lần phân
nhánh xét cả hai giá trị của biến được chọn, phủ hết các mở rộng
chưa bị loại. Khi mỗi mệnh đề đã có ít nhất một literal đúng,
điền các biến chưa gán không làm mất các literal đúng đó.
Vì vậy mô hình trả về thỏa CNF cùng giả định và, nếu mọi nhánh
bị loại, công thức không có nghiệm.

Độ sâu đệ quy không vượt quá $V$, số biến gồm cả biến phụ. Số
nút tìm kiếm xấu nhất tỉ lệ $O(2^V)$. Nếu $L$ là tổng số lần
xuất hiện literal và $M$ là số mệnh đề, mỗi nút có thể quét
$O(V+1)$ lượt với chi phí $O(L+M)$ mỗi lượt, bên cạnh thao tác
sao chép phép gán. Do đó bộ đếm trạng thái và các biến đạt tới
có thể làm tăng chi phí dù số cầu thủ giữ nguyên.

\clearpage
\section{Logic Agent và suy luận từng bước}
\subsection{Chứng minh phán quyết}
Một mô hình SAT chỉ cho biết một cách gán phù hợp với KB. Phán
quyết được chứng minh khi giữ nguyên giá trị trong mọi mô hình.
Agent dựng lại $\KB_t$ bằng \code{build_kb}, kiểm tra ánh xạ chính,
phụ và các literal, rồi bảo đảm mọi phán quyết công khai đều có
mệnh đề đơn vị. Mỗi lần truy vấn, Agent xác minh mô hình SAT trả về
đủ biến và thỏa cả công thức lẫn giả định. Kết quả UNSAT dựa vào
bộ giải và được kiểm tra thông qua các bài thử thuật toán.

Sau khi xác nhận $\KB_t$ là \SAT, \code{_classify} gọi DPLL một lần
với $+i$ và một lần với $-i$, trong đó $i$ là ID biến chính của ô.
Khi $\KB_t$ tự nó là \UNSAT, bước kiểm tra cơ sở trả
\code{INCONSISTENT} trước khi xét hai giả định. Điều này ngăn một
KB mâu thuẫn sinh ra phán quyết trong ván.

\begin{table}[htbp]
\caption{Phân loại một ô chưa giải với KB nhất quán.}
\label{tab:entailment}
\centering
\small
\begin{tabular}{@{}ccll@{}}
\toprule
$\KB_t\land C_i$ & $\KB_t\land\neg C_i$ & \textbf{Phân loại} & \textbf{Kết luận}\\
\midrule
\SAT & \UNSAT & \code{CRIMINAL} & $\KB_t\models C_i$\\
\UNSAT & \SAT & \code{INNOCENT} & $\KB_t\models\neg C_i$\\
\SAT & \SAT & \code{UNKNOWN} & Cả hai trạng thái còn khả dĩ.\\
\bottomrule
\end{tabular}
\end{table}

\subsection{Phán quyết, Hint và Auto Solve}
\begin{enumerate}
\item \textbf{Phán quyết thủ công.} \code{submit} lấy trạng thái
  công khai và gọi \code{classify_cell}. Kết quả có thể là
  \code{ALREADY_SOLVED}, \code{INCONSISTENT}, \code{NOT_PROVABLE},
  \code{CONTRADICTED} hoặc \code{ACCEPTED}. Hai trường hợp cần
  phân biệt khi chơi là cả hai phán quyết còn khả dĩ và phán quyết
  đối lập đã được chứng minh. Trong cả hai trường hợp, thẻ vẫn úp.
\item \textbf{Hint.} \code{hint} kiểm tra KB cơ sở và duyệt theo
  hàng, cột đến ô đầu tiên có phán quyết bắt buộc.
  \code{classify_all} hỗ trợ phân loại toàn bộ các ô chưa giải
  bằng một KB cơ sở chung. Hint có thể làm tăng bộ đếm SAT
  nhưng chưa mở thẻ và chưa ghi một bước chấp nhận.
\item \textbf{Step và Auto Solve.} \code{auto_step} dùng ngay chứng
  minh từ Hint nội bộ để mở tối đa một thẻ. \code{auto_solve}
  lặp các bước trên trạng thái mới cho tới khi hoàn thành hoặc
  dừng. Với cấu hình mặc định và KB nhất quán, nếu ô bắt buộc
  đầu tiên là ô thứ $k$ trong lượt duyệt, bước đó thực hiện
  $1+2k$ truy vấn. Controller có tùy chọn
  \pathcode{check_revealed_consistency} để kiểm tra thêm
  một lần ngay sau khi nhận manh mối mới. Ở cấu hình mặc định,
  KB mới sẽ được kiểm tra ở lần suy luận kế tiếp, nếu có.
\end{enumerate}

Trước khi mở thẻ, Controller so lại snapshot dùng để chứng minh.
Sau khi Engine cập nhật, Controller kiểm tra bàn và dữ kiện cũ
được giữ nguyên, chỉ có một phán quyết và một manh mối mới. Nếu
hợp đồng bị vi phạm sau cập nhật, Controller chặn thao tác tiếp
theo cho tới Restart. Đường xử lý này bảo vệ các lần suy luận
sau và không tự hoàn tác dữ liệu đã bị Engine thay đổi.

\clearpage
\subsection{Nhật ký trên bàn Counting Trail}
Bảng~\ref{tab:trace} trích ba lượt đầu từ lệnh \code{solve}
trên bàn \pathcode{3x3_counting_chain.json}. Mỗi lượt đều bắt đầu
với một truy vấn $\KB_t$ cho kết quả \SAT. Nhật ký
\code{DeductionStep} lưu số bước, ID các manh mối đã có trước
lượt đó, kết quả các truy vấn, ô, phán quyết và ID manh mối vừa mở.
Khi Hint duyệt qua các ô khác, nhật ký của bước cũng lưu các
truy vấn đã thực hiện trong lượt duyệt đó.

\begin{table}[htbp]
\caption{Ba bước đầu trong trace thực tế của 3x3 Counting Trail.}
\label{tab:trace}
\small
\begin{tabularx}{\textwidth}{@{}c L{30mm} Y c c l@{}}
\toprule
\textbf{Bước} & \textbf{Phán quyết} & \textbf{Manh mối hoạt động}
& $\KB_t\land C_i$ & $\KB_t\land\neg C_i$ & \textbf{Mở mới}\\
\midrule
1 & B1 \code{INNOCENT} & \code{clue_A1} & \UNSAT & \SAT & \code{clue_B1}\\
2 & C1 \code{CRIMINAL} & \code{clue_A1}, \code{clue_B1}
& \SAT & \UNSAT & \code{clue_C1}\\
3 & A2 \code{CRIMINAL} & \code{clue_A1}, \code{clue_B1}, \code{clue_C1}
& \SAT & \UNSAT & \code{clue_A2}\\
\bottomrule
\end{tabularx}
\end{table}

Cả ván mở thêm tám thẻ với 24 truy vấn SAT. Số truy vấn này
ứng với một truy vấn cơ sở và hai giả định cho mỗi lượt, vì
ô đầu tiên còn úp luôn chứng minh được trong bài này.

\subsection{Nghiệm duy nhất của bộ manh mối đầy đủ}
\label{sec:uniqueness}
Engine công bố toàn bộ manh mối qua \pathcode{all_clues_for_validation}
trong luồng kiểm định. \code{check_uniqueness} mã hóa các manh
mối đó với metadata bàn và bỏ phán quyết hiện tại của ván.
Nếu CNF đầy đủ là \UNSAT, bàn không có nghiệm.
Nếu SAT với mô hình $M$, truy vấn thứ hai thêm một mệnh đề chặn
phép gán đó trên \emph{toàn bộ biến chính}
\[
 B_M=\bigvee_{i=1}^{N^2}
 \begin{cases}
  \neg C_i & \text{khi } M(C_i)=1,\\
  C_i & \text{khi } M(C_i)=0.
 \end{cases}
\]
Truy vấn thứ hai UNSAT xác nhận nghiệm duy nhất. Kết quả SAT
cung cấp một lời giải khác trên các ô. Biến phụ được bỏ qua khi
chặn vì khác biệt ở biến phụ không làm đổi lời giải của bàn.
Để tính nhãn mở đầu như tiên đề của bài kiểm định đầy đủ, các
nhãn ấy cần được biểu diễn trong chính tập manh mối được mã hóa.

Tính duy nhất trên toàn bộ manh mối và khả năng giải từng bước là
hai phép kiểm định riêng. Ngay cả khi công thức đầy đủ chỉ có
một nghiệm, ở một thời điểm của ván có thể vẫn chưa có ô nào
được ép bởi $\KB_t$. Khi đó Agent dừng và giữ thẻ úp.
Lệnh \code{validate} kiểm tra cả nghiệm duy nhất lẫn khả năng
hoàn thành bằng mở thẻ tuần tự.

\subsection{Phạm vi cách ly dữ liệu}
Engine giữ lời giải và manh mối chưa mở. Agent nhận
\code{PublicGameState} trong \code{classify_cell}, \code{hint}
và \code{classify_all}, rồi sao chép các container công khai
trước khi chuyển cho Encoder. Controller kiểm tra bằng Agent
trước khi Engine dùng lời giải ẩn làm hậu kiểm. Nút kiểm tra
nghiệm duy nhất dùng API quản trị tách khỏi luồng chơi. API đó
có thể cập nhật thống kê truy vấn nhưng không đưa mô hình
của toàn bộ manh mối vào Hint hay mở thêm thẻ. Ranh giới này
được thực thi bằng dữ liệu truyền qua các giao diện trong một
tiến trình Python.

\clearpage
\section{Game Engine và giao diện trình duyệt}
\label{sec:web}
\code{main.py} mặc định mở bàn \emph{4x4 Rookie Training}.
\code{run_gui} khởi động máy chủ HTTP cục bộ, in URL rồi mở trình
duyệt. Cổng ưu tiên là 8765. Khi cổng đã được dùng, máy chủ chọn
một cổng trống. Phần giao diện gồm HTML, CSS và JavaScript tĩnh
cùng tài nguyên ảnh trong \pathcode{assets/}. Danh sách 25 cầu thủ
và các đội hình cố định được định nghĩa trong \pathcode{roster.py}.
Thẻ có ảnh nếu tài nguyên cục bộ đọc được, hoặc hiển thị áo đấu
SVG dự phòng. Tên cầu thủ là dữ liệu trình bày, còn nhãn
\code{CRIMINAL} và \code{INNOCENT} thuộc ván chơi hư cấu~\cite{credits}.

\begin{table}[htbp]
\caption{Thao tác trên giao diện và xử lý phía máy chủ.}
\label{tab:actions}
\small
\begin{tabularx}{\textwidth}{@{}L{35mm}Y@{}}
\toprule
\textbf{Thao tác} & \textbf{Thực hiện trên bàn chơi}\\
\midrule
Load, Restart & Load chọn một JSON trong danh mục và tạo Engine, Agent
cùng Controller mới. Restart đưa bàn hiện tại về các thẻ mở đầu,
xóa trace và đặt lại thống kê.\\
Chơi thủ công & Chọn trên sân hoặc danh sách rồi gửi ô và phán quyết.
Khi \code{ACCEPTED}, thẻ đổi trạng thái và lật để hiện manh mối.
Thông báo phân biệt nước chưa chứng minh với nước bị bác bỏ.\\
Chọn manh mối & \code{Match} tuần tự hóa văn bản, công thức, ô
tham chiếu và các vùng. JavaScript dùng dữ liệu công khai để tô sáng
các ô và hiển thị đường nối trên sân.\\
Hint, Step & Hint làm nổi ô có phán quyết bắt buộc và các thẻ đang
cung cấp dữ kiện. Step chứng minh và mở tối đa một thẻ.\\
Auto Solve, Instant & Auto Solve gửi từng yêu cầu \code{step} theo
nhịp, có thể dừng giữa hai lượt. Instant gửi \code{solve} để chạy
vòng giải trong một yêu cầu.\\
Trace, KB, nghiệm duy nhất & Các tab hiển thị lượt mở, truy vấn,
số đo và tối đa 30 mệnh đề đầu của KB công khai. Nút Check uniqueness
mở hộp kết quả từ phép kiểm định toàn bộ manh mối.\\
\bottomrule
\end{tabularx}
\end{table}

Trình duyệt đọc \pathcode{/api/puzzles} và \pathcode{/api/state},
sau đó gửi JSON cho các thao tác và dựng lại giao diện từ snapshot
trả về. \code{Match} giữ một ván trong tiến trình. Các thao tác
được tuần tự hóa bằng khóa. Server trả lỗi JSON cho đầu vào không
hợp lệ. Khi phát hiện ID phiên cũ, server từ chối thao tác với
HTTP 409 và giao diện đồng bộ lại. Khi mất kết nối, giao diện dừng
Auto Solve, báo lỗi và yêu cầu tải lại trạng thái trước lượt chơi
tiếp theo.

\paragraph{Phạm vi xác minh.}
Pytest kiểm tra xử lý yêu cầu HTTP, trạng thái công khai, thao
tác chơi, phiên cũ và ảnh cầu thủ. Các chi tiết về hiệu ứng lật
thẻ, tô sáng và bàn 2.5D ở đây dựa trên mã giao diện. Đợt rà soát
này chưa xác nhận trực tiếp thao tác bằng một phiên trình duyệt.

\clearpage
\section{Thực nghiệm và kiểm thử}
\label{sec:experiments}
\subsection{Thiết lập và kết quả đo trên phiên bản hiện tại}
Tập so sánh cố định trong \pathcode{experiments/puzzle_set.json}
có hai bàn $3\times3$, hai bàn $4\times4$ và một bàn $5\times5$.
Đo lại ngày 26/11/2026 trên commit \code{b373f86}, Python 3.14.5,
Windows 11, với năm lượt cho mỗi bàn và giới hạn 60 giây mỗi lượt.
Bộ đếm và thời gian được đo từ luồng \code{auto_solve} mặc định.
Các bảng dùng mã ở Bảng~\ref{tab:cases} để định danh tệp đầu vào.

\begin{table}[htbp]
\caption{Tập bàn và cấu trúc CNF của toàn bộ manh mối.}
\label{tab:cases}
\centering\small
\begin{tabularx}{\textwidth}{@{}lYcrrr@{}}
\toprule
\textbf{Mã} & \textbf{Tệp JSON} & \textbf{Cỡ}
& \textbf{Biến chính} & \textbf{Biến phụ} & \textbf{Mệnh đề}\\
\midrule
B3 & \pathcode{3x3_binary_chain.json}    & $3\times3$ & 9  & 0     & 16\\
C3 & \pathcode{3x3_counting_chain.json}  & $3\times3$ & 9  & 66    & 185\\
R4 & \pathcode{4x4_easy_training.json}   & $4\times4$ & 16 & 315   & 1\,611\\
L4 & \pathcode{4x4_medium_links.json}    & $4\times4$ & 16 & 850   & 5\,460\\
D5 & \pathcode{5x5_deep_inference.json}  & $5\times5$ & 25 & 1\,539 & 13\,969\\
\bottomrule
\end{tabularx}
\end{table}

\begin{table}[htbp]
\caption{Số đếm khi giải từ manh mối công khai theo từng lượt.}
\label{tab:work}
\centering\small
\begin{tabular}{@{}lrrrrrrc@{}}
\toprule
\textbf{Mã} & \textbf{SAT calls} & \textbf{Decisions} & \textbf{Propag.} &
\textbf{Conflicts} & \textbf{Backtracks} & \textbf{Bước} & \textbf{Lượt xong}\\
\midrule
B3 & 24  & 0      & 132      & 8      & 0      & 8  & 5/5\\
C3 & 24  & 0      & 924      & 8      & 0      & 8  & 5/5\\
R4 & 93  & 222    & 8\,578   & 33     & 18     & 15 & 5/5\\
L4 & 153 & 863    & 41\,377  & 174    & 297    & 15 & 5/5\\
D5 & 268 & 5\,683 & 183\,366 & 3\,832 & 7\,088 & 24 & 5/5\\
\bottomrule
\end{tabular}
\end{table}

\begin{table}[htbp]
\caption{Thời gian trung bình $\pm$ độ lệch chuẩn tổng thể (ms).}
\label{tab:runtime}
\centering\small
\begin{tabular}{@{}lrrl@{}}
\toprule
\textbf{Mã} & \textbf{Tổng thời gian DPLL} & \textbf{Vòng giải tự động} & \textbf{Kết thúc}\\
\midrule
B3 & $0{,}4544 \pm 0{,}0213$ & $2{,}8680 \pm 0{,}1740$ & \code{SOLVED}\\
C3 & $2{,}0420 \pm 0{,}1743$ & $6{,}5186 \pm 0{,}4576$ & \code{SOLVED}\\
R4 & $45{,}5838 \pm 2{,}7301$ & $87{,}2889 \pm 7{,}3876$ & \code{SOLVED}\\
L4 & $434{,}7492 \pm 10{,}3870$ & $649{,}1232 \pm 19{,}8064$ & \code{SOLVED}\\
D5 & $5\,249{,}2469 \pm 114{,}8651$ & $6\,816{,}4065 \pm 146{,}3456$ & \code{SOLVED}\\
\bottomrule
\end{tabular}
\end{table}

Số biến và mệnh đề ở Bảng~\ref{tab:cases} mô tả công thức toàn bộ
manh mối. Các số đếm ở Bảng~\ref{tab:work} mô tả truy vấn trên KB
công khai, và runner lấy số đếm từ lượt hoàn tất đầu tiên.
Thời gian DPLL là tổng thời gian các lệnh \code{solve}. Thời gian
vòng giải còn bao gồm lập KB, kiểm tra mô hình và mở thẻ, rồi
dừng đồng hồ khi vòng giải kết thúc. Khâu nạp và đo kích thước
CNF đầy đủ nằm ngoài cả hai số thời gian. Trong 25 lượt đo
không xảy ra lỗi, phiên dừng \code{UNKNOWN} hoặc timeout.

\clearpage
\subsection{Đối chiếu và giải thích chênh lệch}
Hai bàn $3\times3$ có cùng tám lượt mở mới và 24 truy vấn SAT.
B3 dùng chủ yếu các ràng buộc nhị phân, còn C3 dùng bộ đếm nên
số lan truyền tăng từ 132 lên 924. Cả hai bàn đều được giải bằng
lan truyền trong mọi truy vấn đã thống kê, không phát sinh
quyết định rẽ nhánh.

R4 và L4 có cùng 16 biến chính và 15 lượt mở mới nhưng L4 cần
153 truy vấn, so với 93 của R4. Agent phải xét thêm các ô chưa
chứng minh được trước khi tìm thấy một bước bắt buộc. L4 có
nhiều biến phụ, nhiều mệnh đề hơn và gần bốn lần số quyết định.
Bàn D5 có các vùng và manh mối liên thông, đếm và so sánh.
Nó tạo công thức lớn nhất và có số quay lui cao nhất. Các số
đếm cho thấy chi phí tìm kiếm phụ thuộc vào cấu trúc ràng buộc
cùng thứ tự tìm nước đi, bên cạnh kích thước bàn.

Ba tệp trong \pathcode{results/} lưu cùng đợt benchmark ngày
20/11/2026, phát sinh trước bản tích hợp hiện tại~\cite{results}.
Đợt đó có 681 biến phụ và 4\,838 mệnh đề ở L4, 1\,442 biến
phụ và 13\,493 mệnh đề ở D5. Ở phiên bản hiện tại, cùng hai
bàn tạo 850 và 1\,539 biến phụ. \code{COUNT_COMPARE} đã chuyển
từ liệt kê các phép gán vùng sang hai bộ đếm cùng các mệnh đề
loại cặp số đếm, làm thay đổi cấu trúc CNF. Các thời gian từ hai
đợt thuộc hai phiên bản mã và môi trường đo, nên chúng chỉ mô
tả từng lần đo.

\subsection{Kiểm thử tự động và ca chẩn đoán}
\label{sec:tests}
\code{pytest} 8.4.2 thu được \textbf{119 test đạt, 0 thất bại,
0 bỏ qua} trong 25,74 giây. Phạm vi gồm phân giải vùng, bảng chân
trị đối chiếu CNF, SAT/UNSAT và giả định, Agent/Controller, cách ly
dữ liệu, hợp đồng mở thẻ, CLI, ghi báo cáo benchmark và xử lý Web API.
Suite có ca giả lập timeout để kiểm tra runner giữ ca lỗi trong bảng.
\code{python main.py validate puzzles} xác nhận cả năm bàn chuẩn
đều thỏa được, có nghiệm duy nhất và giải tuần tự được với số bước
8, 8, 15, 15, 24 theo thứ tự của Bảng~\ref{tab:cases}.

\begin{table}[htbp]
\caption{Hai ca chẩn đoán nằm ngoài tập benchmark cố định.}
\small
\begin{tabularx}{\textwidth}{@{}L{38mm}Y Y@{}}
\toprule
\textbf{Ca} & \textbf{Kết quả CLI} & \textbf{Kết quả phía Web API}\\
\midrule
\emph{Ambiguous VAR check}
& \code{validate} báo nhiều nghiệm và không thể hoàn tất
tuần tự. \code{solve} dừng \code{UNKNOWN}, 17 lần gọi SAT, 0 bước.
& Hint báo \code{NOT_PROVABLE}. Kiểm định đầy đủ trả
\code{AMBIGUOUS}. Số thẻ mở giữ ở 1/9.\\
\emph{Inconsistent match reports}
& Chế độ nạp thông thường từ chối vì \code{clue_B1} sai theo
lời giải lưu trong JSON.
& \code{Match} cho phép mở ở chế độ chẩn đoán riêng. Hint và
kiểm định đầy đủ báo \code{INCONSISTENT}, số thẻ mở giữ ở 2/9.\\
\bottomrule
\end{tabularx}
\end{table}

Các kết quả chẩn đoán minh họa hai điểm khác nhau của vòng suy
luận. KB có nghiệm nhưng thiếu kết luận buộc sẽ dừng mà giữ kín
các thẻ còn lại. KB mâu thuẫn cần được báo là lỗi nhất quán trước
khi suy diễn bất cứ phán quyết nào.

\clearpage
\section{Kết luận và giới hạn}
Ở phiên bản được khảo sát, Engine, Controller, Logic Agent, Encoder
và DPLL tích hợp thành một vòng chơi không đoán. Hệ thống phân biệt
nước đi thiếu chứng cứ, nước đi trái với kết luận và KB mâu thuẫn.
Các bàn chuẩn đều vượt qua cả kiểm định nghiệm duy nhất lẫn kiểm
định mở thẻ từng bước. Các ca chẩn đoán cho thấy chương trình
có thể dừng khi chứng cứ không đủ hoặc khi manh mối mâu thuẫn.

Cài đặt xây lại CNF cho từng lượt và quét toàn bộ mệnh đề trong
mỗi lần lan truyền. Bộ đếm sinh mệnh đề loại trừ theo từng cặp,
còn liên thông thêm biến đạt tới theo độ sâu. Cây tìm kiếm DPLL
vẫn có thể tăng theo hàm mũ. Thực nghiệm chỉ gồm năm bàn được
chọn, với $5\times5$ là cỡ lớn nhất đã đo. Kết quả định lượng
áp dụng cho tập đó và cấu hình đã nêu. Instant giữ yêu cầu HTTP
đang chạy cho tới khi vòng giải kết thúc.

Trên các bàn tùy chỉnh, phần tô sáng của \code{GROUP_MOST} và
\code{GROUP_ONLY} lấy nhóm được nêu trong manh mối. Các nhóm so
sánh ngầm trong công thức không được đánh dấu đầy đủ. Phần kiểm thử
Web API cũng chưa đo khả năng sử dụng hoặc độ tương thích giữa
các trình duyệt. Để mở rộng quy mô, hướng khả thi là giữ lại phần
CNF đã lập theo ván, thử cách mã hóa đếm gọn hơn và giảm số lượt
quét bằng literal theo dõi. Một tập kiểm thử lớn hơn theo nhóm
manh mối sẽ giúp phân biệt ảnh hưởng của cỡ bàn và độ phức tạp
ràng buộc.

\phantomsection
\addcontentsline{toc}{section}{Tài liệu tham khảo}
\begin{thebibliography}{9}
\bibitem{original}
Knotwise Games. \emph{Griductive: How to Play}.
\href{https://knotwise.games/games/griductive/how-to-play}{Hướng dẫn trò chơi gốc}.

\bibitem{spec}
Khoa Công nghệ Thông tin, Trường Đại học Khoa học Tự nhiên,
ĐHQG-HCM. \emph{CSC14003 | Introduction to Artificial Intelligence.
Project 2: Griductive Solver}. Tệp đề bài
\code{AI - Project 2.pdf}, các mục 2--4.7.

\bibitem{dpll}
M. Davis, G. Logemann và D. W. Loveland.
\emph{A machine program for theorem-proving}.
Communications of the ACM, 5(7), 394--397, 1962.
\href{https://doi.org/10.1145/368273.368557}{DOI 10.1145/368273.368557}.

\bibitem{aima}
S. J. Russell và P. Norvig.
\emph{Artificial Intelligence: A Modern Approach}, ấn bản thứ 4,
Pearson, chương 7.
\href{https://aima.cs.berkeley.edu/}{Trang chính thức của sách}.

\bibitem{results}
Nhóm Griductive Solver. Mã nguồn và bộ kiểm thử tại commit
\code{b373f86} của \href{https://github.com/kitchauhcmus/griductive-solver}
{kho mã dự án}. Dữ liệu benchmark ngày 20/11/2026 trong
\pathcode{results/benchmark_results.csv},
\pathcode{results/benchmark_results.json} và
\pathcode{results/EXPERIMENTS.md}.

\bibitem{credits}
Nhóm Griductive Solver. Tác giả ảnh, nguồn và giấy phép của tài nguyên
cầu thủ và sân vận động trong \pathcode{assets/CREDITS.md}.
\end{thebibliography}

\clearpage
\appendix
\section{Sử dụng AI tạo sinh, tái lập và video}
\subsection{Phạm vi hỗ trợ và nội dung yêu cầu}
Trong lần cập nhật báo cáo này, nhóm dùng Codex để đối chiếu
mã nguồn với đề bài, hỗ trợ biên tập tiếng Việt và mã LaTeX,
chạy lại kiểm thử và benchmark, rồi biên dịch và kiểm tra PDF.
Các số liệu được lấy từ đầu ra của chương trình. Phân công và
MSSV được giữ theo hồ sơ nhóm. Tỷ lệ hoàn thành từng thành viên
chờ xác nhận từ nhóm.

Các câu sau trích từ yêu cầu của lần rà soát này.
\begin{quote}
Hãy phân tích toàn bộ mã nguồn thực tế trong workspace hiện tại.\\
Mã nguồn hiện tại là nguồn sự thật chính.\\
Chỉ được chỉnh sửa và cập nhật hai file [\texttt{Report.tex} và \texttt{Report.pdf}].
\end{quote}
Tên hai tệp trong câu cuối được rút gọn từ một danh sách.
Lịch sử sử dụng AI trong giai đoạn phát triển cần nhóm xác nhận.

\subsection{Tái lập kết quả}
Chạy từ thư mục gốc của dự án. Lệnh benchmark dưới đây in bảng
JSON ra terminal. Thêm các tùy chọn \code{--csv}, \code{--json}
và \code{--markdown} với đường dẫn đầu ra riêng để lưu các dạng
báo cáo của một lần chạy.
\begingroup
\footnotesize
\begin{verbatim}
python -B -m pytest -q -p no:cacheprovider
python -B main.py validate puzzles
python -B main.py solve puzzles/3x3_counting_chain.json
python -B main.py experiments --manifest experiments/puzzle_set.json --runs 5 --timeout 60
\end{verbatim}
\endgroup

\subsection{Video minh họa}
\textbf{Liên kết Google Drive của video demo} \emph{Chờ nhóm bổ sung.}

\end{document}
